\section{Experimental Setup}
\label{sec:setup}

\subsection{States and District Counts}

We evaluate all algorithms on five states spanning the realistic range of
congressional district counts and two distinct geographic characters:

\begin{itemize}
  \item \textbf{North Carolina (NC)}, $k = 14$.
        Baseline contested state from G.14.
        $N \approx 2{,}200$ census tracts; competitive partisan baseline.
        G.14 search results carried forward; structure and new search
        results are new.

  \item \textbf{Wisconsin (WI)}, $k = 8$.
        Small-$k$ baseline from G.14.
        $N \approx 1{,}500$ census tracts.
        Tests overhead cost of algorithms designed for larger $k$.
        G.14 search results carried forward.

  \item \textbf{Texas (TX)}, $k = 38$.
        Large-$k$ canonical test case from G.14.
        $N \approx 5{,}200$ census tracts.
        Primary state for autocorrelation and Parallel Tempering comparisons.
        G.14 search results carried forward; new algorithms added.

  \item \textbf{Florida (FL)}, $k = 28$.~\dag{}
        Intermediate-$k$ state testing generalisability of TX findings.
        $N \approx 4{,}100$ census tracts; complex coastal geometry
        affects \pp{} for all algorithms.
        All results are new.

  \item \textbf{New Hampshire (NH)}, $k = 2$.~\dag{}
        The ILP-tractable state.
        $N \approx 260$ census tracts; each bisection half has $\leq 130$
        tracts, well within the ILP tractability threshold of $N \leq 500$.
        NH is the only state for which certified-optimal \ec{} is available.
        All results are new; ILP comparison is the primary contribution
        for this state.
\end{itemize}

\dag{} denotes new single-run estimates not present in G.14.

\subsection{Experimental Protocol}

\paragraph{Structure evaluation.}
For the structure-layer comparison (Section~\ref{sec:structure}), we run each
structure algorithm with the \texttt{--search convergence} flag (the
ConvergenceSweep federal statute default) at fixed convergence threshold
$T_{\mathrm{conv}} = 600$.
This ensures that the search layer does not confound the structure comparison;
all plans reach the same convergence criterion before metrics are recorded.
Exception: ILP is run with \texttt{--search single} (NH only), since
convergence sweep on a single-node ILP plan is unnecessary.

\paragraph{Search evaluation.}
For the search-layer comparison (Section~\ref{sec:search}), we fix
\texttt{--structure prime-factor} (ApportionRegions bisection) and
\texttt{--weights county} across all search runs.
This follows the G.14 protocol (which used \texttt{--structure ratio-optimal}
and \texttt{--weights geographic}) updated to the ApportionRegions composite
which has become the platform default for production runs.
G.14 results were recorded under the original protocol and are carried forward
without adjustment; the directional comparisons are robust to this protocol
difference as established by the G.14 seed-robustness analysis.

\textit{Directional robustness:} G.14's seed-stability analysis (5 seeds,
NC/WI/TX) found that EC and PP orderings are stable across seeds with
${<}5\%$ variation; this gives confidence that the algorithm-level comparisons
carried from G.14 are directionally correct even under the protocol change.
Full re-runs under the G.15 protocol are deferred to the Phase~2 empirical
study.

\paragraph{Common configuration.}
\begin{itemize}
  \item \textbf{Base seed}: 42 (all new runs).
  \item \textbf{Population tolerance}: $\epsilon = 0.5\%$.
  \item \textbf{Hardware}: Intel Core i7-11800H (8 cores, 16 threads);
        structure runs use \texttt{--workers 1} (single-threaded);
        search runs use \texttt{--workers 16} (full Rayon parallelism)
        for \pt{} and \ams{} which benefit from multi-core; all other
        search runs use \texttt{--workers 1} to match G.14.
  \item \textbf{Census year}: 2020.
  \item \textbf{Chain budget}: $T = 1{,}000$ steps for all chain methods.
  \item \textbf{SMC particle budget}: $N = 1{,}000$ particles.
  \item \textbf{PT replicas}: $R = 8$ replicas across temperature ladder
        $\beta \in \{0.1, 0.2, 0.35, 0.5, 0.65, 0.8, 0.9, 1.0\}$.
\end{itemize}

\subsection{Metrics}

Five metrics are reported; the first four appeared in G.14.

\paragraph{Edge-cut score (EC).}
$\ec(\pi) = |\{(u,v) \in E : \pi(u) \neq \pi(v)\}|$.
Lower \ec{} is more compact (fewer cross-district adjacencies).

\paragraph{Mean Polsby-Popper (PP).}
$\pp = \frac{1}{k} \sum_{d=1}^k \frac{4\pi A_d}{P_d^2}$.
Higher \pp{} is more compact (closer to a circle).

\paragraph{Democratic seat count (D-seat).}
Democratic congressional seat count under 2020 presidential vote totals,
applied to the final plan of each algorithm.
Reported as a raw count (not stability fraction) to show partisan sensitivity
to algorithm choice.

\paragraph{Lag-100 autocorrelation (ACF(100)).}
$\acf(100) = \frac{\sum_{t=1}^{T-100}(\ec_t - \bar{\ec})(\ec_{t+100}
- \bar{\ec})}{\sum_{t=1}^T (\ec_t - \bar{\ec})^2}$.
Lower is better (faster mixing). Not applicable to SMC or structure algorithms.

\paragraph{Wall-clock runtime (seconds).}
Total wall clock time from job start to final plan output, on the hardware
specified above.
For \pt{}, runtime is measured with \texttt{--workers 16} since single-threaded
PT is not representative of production use.

\subsection{Algorithm--CLI Mapping}

Table~\ref{tab:cli} gives the CLI flags for all 15 algorithms.

\begin{table}[ht]
\centering
\caption{CLI flags for all 15 algorithms.}
\label{tab:cli}
\small
\begin{tabular}{@{}llll@{}}
\toprule
Algorithm & Layer & Flag & Paper \\
\midrule
Standard METIS       & Structure & \texttt{--structure standard-bisect}    & (default) \\
Simulated Annealing  & Structure & \texttt{--structure sa-bisect}           & B.19 \\
CVD Graph-Distance   & Structure & \texttt{--structure cvd-graph}           & B.22 \\
CVD Geographic       & Structure & \texttt{--structure cvd-geographic}      & B.22b \\
BFS Growth           & Structure & \texttt{--structure bfs-growth}          & B.23 \\
ILP Exact            & Structure & \texttt{--structure ilp}                 & B.24 \\
\midrule
ConvergenceSweep     & Search & \texttt{--search convergence}           & B.16 \\
PercentileSweep      & Search & \texttt{--search percentile}            & H.0 \\
BisectionEnsemble    & Search & \texttt{--search bisection-ensemble}    & H.1 \\
Short-Burst          & Search & \texttt{--search short-burst}           & G.6 \\
Flip                 & Search & \texttt{--search flip}                  & G.8 \\
Forest ReCom         & Search & \texttt{--search forest-recom}          & G.9 \\
Merge-Split          & Search & \texttt{--search merge-split}           & G.10 \\
Multi-scale          & Search & \texttt{--search multiscale}            & G.11 \\
ShortBurstForest     & Search & \texttt{--search shortburst-forest}     & G.12 \\
VRA-Aware ReCom      & Search & \texttt{--search vra-recom}             & G.13 \\
Parallel Tempering   & Search & \texttt{--search parallel-tempering}    & B.20 \\
SMC-Percentile       & Search & \texttt{--search smc-percentile}        & G.7/new \\
Adaptive Multi-scale & Search & \texttt{--search adaptive-multiscale}   & B.21 \\
\bottomrule
\end{tabular}
\end{table}

\subsection{Single-Run Caveat}
\label{sec:caveat}

All results in this paper are single-run (seed 42) and should be treated as
illustrative.
Results labelled \dag{} are new estimates for FL, NH, or new search algorithms;
results without \dag{} are carried from G.14 where they were cross-validated
across seeds 42--46.
The decision framework in Section~\ref{sec:decision} is grounded in the
structural properties of each algorithm (established in their respective
B/G/H-series papers), not solely in the single-run empirics here.
